% !TeX root = courtade-kumar.tex
\section{Preliminaries}
\label{sec:preliminaries}
We use $\ent{\cdot}$ and $\muti{\cdot}{\cdot}$ for Shannon
entropy and mutual information, with logarithms in base $2$.
We denote the natural logarithm by $\ln$. We only need
Tsallis entropy for bits, while the observed variables may
be arbitrary. For a bit $C$ with mean $\mu$, we define
\begin{equation*}
  \tsent{C}=4\mu(1-\mu).
\end{equation*}
For a bit $C$ and an arbitrary random variable $Z$,
Tsallis information is defined by
\begin{equation*}
  \tsmuti{C}{Z}=\tsent{C}-\tsent{C\emid Z}.
\end{equation*}
Note that for bits $B$, $C$, the quantities
$\tsmuti{B}{C}$ and $\tsmuti{C}{B}$ need
not be equal. We denote
\begin{equation*}
\begin{aligned}
  \sI=1-\entb{p},\qquad \qI&=(1-2p)^2,\\
  \text{and}\qquad \qJ&=\qI-\sI.
\end{aligned}
\end{equation*}
Thus $\sI=\muti{Y_i}{X_i}$ and $\qI=\tsmuti{Y_i}{X_i}$.
With these shorthands, we have
\begin{equation*}
  \qent{C}=\qI\ent{C}-\qJ\tsent{C}.
\end{equation*}
For all our entropy and information quantities
$\ent{\cdot}$, $\muti{\cdot}{\cdot}$, $\tsent{\cdot}$,
$\tsmuti{\cdot}{\cdot}$, and $Q$, conditioning on $U$ means
evaluating the quantity under the conditional law given $U$
and taking expectation over $U$, with $p$ held fixed.
For instance,
$\qent{C\emid U} = \E_{u\sim U}[\qent{C\emid U=u}]$.
Next, we gather some facts about Tsallis and Shannon
information.
\begin{lemma}
\label{lem:tsallis-properties}
Let $B$ and $C$ be bits and $Z$, $W$, $U$ be arbitrary
random variables. We have
\leavevmode
\begin{enumerate}
\setlength{\itemsep}{4pt}
\item\label{item:tsallis-bounds}
$0\le\tsmuti{C}{Z\emid U}\le\tsent{C\emid U}$ and
$\tsent{C\emid U}\le\ent{C\emid U}
\le\bigl(1+\tsent{C\emid U}\bigr)/2$.
\item\label{item:tsallis-chain}
$\tsmuti{C}{Z,W\emid U}
=\tsmuti{C}{Z\emid U}+\tsmuti{C}{W\emid Z,U}$.
If $C-Z-W$ is a Markov chain, then
$\tsmuti{C}{W}\le\tsmuti{C}{Z}$.
\item\label{item:tsallis-shannon}
$\tsmuti{C}{Z\emid U}\le2\ln2\,\muti{C}{Z\emid U}$.
If $C$ is unbiased and independent of $U$, then also
$\muti{C}{Z\emid U}\le\tsmuti{C}{Z\emid U}$.
\item\label{item:tsallis-bits}
$\tsent{C}\tsmuti{B}{C}=\tsent{B}\tsmuti{C}{B}$.
If $B$ is unbiased, then
$\tsmuti{C}{B}=\bigl(2\Pr[B=C]-1\bigr)^2$.
\item\label{item:tsallis-conditioning}
If $B$ is independent of $U$, then
$\tsmuti{B}{C\emid U}\,\tsent{C\emid U}
\ge\tsent{B}\,\tsmuti{C}{B}$.
\item\label{item:tsallis-noise}
$0\le\qJ\le\sI$ for every $p\in[0,1]$.
\end{enumerate}
\end{lemma}

Proofs of the preceding and the next lemma are given in
\autoref{sec:quadi}. Let us define
\begin{equation*}
  c(\mu)=\frac{\ln(1/\mu)}{4(1-\mu)^2}
  \qquad\nparen{1/2\le\mu<1},
\end{equation*}
extending it to $0<\mu<1/2$ by $c(\mu)=c(1-\mu)$.

\begin{lemma}
\label{lem:biased-comparison}
For a nonconstant bit $C$ with mean $\mu$ and every random variable $U$,
\begin{equation*}
  \muti{C}{U}\le\frac{c(\mu)}{\ln2}\tsmuti{C}{U},
  \qquad\text{and}\qquad
  \tsent{C}\,c(\mu)\le1.
\end{equation*}
\end{lemma}
